\section{Estructura cardinal del continuo y no fidelidad de la reducci\'on extensional}
\label{sec:estatuto-cardinal-zfc}

\subsection{Construcci\'on total del portador y realizaci\'on arquimediana}

Sea \(\mathcal A_{\rm HMT}\) el atlas graduado de memorias Hensel completas
producido por el transductor integral de la arquitectura APP--TRIT--TPK y sea
\begin{equation}
 P_{\rm HMT}:=P_{\rm arq}:
 \mathcal A_{\rm HMT}\longrightarrow\mathbb R_{\rm HMT}
 \label{eq:evaluacion-hmt-cardinal}
\end{equation}
la evaluaci\'on arquimediana obtenida por intersecci\'on de cilindros
encajados. La relaci\'on de equivalencia inducida por la evaluaci\'on y su
cociente son
\begin{equation}
 x\sim_{P_{\rm HMT}}y
 \Longleftrightarrow P_{\rm HMT}(x)=P_{\rm HMT}(y),
 \qquad
 Q_{\rm HMT}:=\mathcal A_{\rm HMT}/\!\sim_{P_{\rm HMT}}.
 \label{eq:quotient-enriquecido-cardinal-section}
\end{equation}
El \cref{thm:cociente-arquimediano-atlas-nuclear} prueba, mediante la
construcci\'on interna precedente, que \(P_{\rm HMT}\) es continua, sobreyectiva y
cociente. Por consiguiente,
\begin{equation}
 \boxed{
 Q_{\rm HMT}\xrightarrow{\ \overline P_{\rm HMT}\ }\mathbb R_{\rm HMT}
 \xrightarrow{\ \jmath_{\rm arq}\ }\mathbb R
 }
 \label{eq:cardinal-cociente-hmt}
\end{equation}
es una sucesi\'on de isomorfismos de espacios ordenados completos; el segundo
morfismo es el reconocimiento arquimediano posterior del cuerpo construido.
La recta \(\mathbb R\) es, por tanto, la imagen extensional del cociente de
memorias Hensel brutas producido por el transductor integral
APP--TRIT--TPK.

La cardinalidad del cociente se obtiene de la propia construcci\'on. En cada
carta,
\begin{equation}
 X^{\rm cal}\cong\mathsf T^{\mathbb N},
 \qquad |\mathsf T|=3,
 \qquad
 \mathcal A_{\rm HMT}
 =\coprod_{m\in\mathbb Z_{\rm HMT}}X^{\rm cal}.
 \label{eq:estructura-cardinal-atlas-hmt}
\end{equation}
Como \(\mathbb Z_{\rm HMT}\) es numerable,
\begin{equation}
 |\mathcal A_{\rm HMT}|
 =\aleph_0\,3^{\aleph_0}=2^{\aleph_0}.
 \label{eq:cardinal-ramas-compatibles}
\end{equation}
Las dos expansiones ternarias de cada extremo racional forman una familia
numerable de pares; su identificaci\'on no altera el cardinal. De
\eqref{eq:cardinal-cociente-hmt} se deduce
\begin{equation}
 \boxed{|Q_{\rm HMT}|=2^{\aleph_0}.}
 \label{eq:cardinal-cociente-hmt-valor}
\end{equation}
Esta igualdad cuenta historias completas m\'odulo igualdad de valor, mientras
que
\begin{equation}
 \left|\bigcup_{n\geq1}\mathsf T^n\right|=\aleph_0,
 \label{eq:prefijos-finitos-numerables}
\end{equation}
cuenta exclusivamente prefijos finitos. La ley generadora puede tener
una presentaci\'on finita y una coalgebra no numerable de prolongaciones
completas.

La formulaci\'on extensional de la hip\'otesis del continuo sobre el portador
realizado se expresa ya en ZF como
\begin{equation}
 \mathsf{CH}_{\rm ext}
 \quad\Longleftrightarrow\quad
 2^{\aleph_0}=\aleph_1
 \quad\Longleftrightarrow\quad
 \nexists A\in\mathcal P(Q_{\rm HMT})\ 
 \bigl(\aleph_0<|A|<|Q_{\rm HMT}|\bigr).
 \label{eq:ch-exacta-sobre-portador-hmt}
\end{equation}
La ecuaci\'on \eqref{eq:ch-exacta-sobre-portador-hmt} se formula despu\'es de
construir el portador y cuantifica sobre la totalidad extensional de sus
subconjuntos una vez olvidada la genealog\'ia. Las equivalencias posteriores
entre igualdad cardinal, inyecci\'on en \(\omega_1\) y buen orden utilizan el
axioma de elecci\'on y se formulan, por tanto, en ZFC.

\begin{theorem}[Separaci\'on de tipos entre generaci\'on y ordinalizaci\'on]
\label{thm:separacion-construccion-ordinalizacion}
La construcci\'on de \(\mathfrak C_{\rm disc}\), la evaluaci\'on total
\(P_{\rm HMT}\), el cociente \(Q_{\rm HMT}\cong\mathbb R_{\rm HMT}\) y
los productos dimensionales
\(Q_{\rm HMT}^{\,d}\cong\mathbb R_{\rm HMT}^{\,d}\) no presuponen
\(\omega_1\), un buen orden del continuo ni la igualdad
\(2^{\aleph_0}=\aleph_1\).
\end{theorem}

\begin{proof}
La estructura discreta conjunta procede del l\'imite inverso de los estados
enriquecidos finitos. La evaluaci\'on procede de la contracci\'on de
cilindros racionales, y el cociente, de la igualdad de sus intersecciones.
Los productos tipados conservan estas operaciones componente a componente.
Ninguna de esas construcciones contiene un ordinal no numerable. La
ordinalizaci\'on pertenece exclusivamente a la representaci\'on extensional
posterior.
\end{proof}

\subsection{Subconjuntos geneal\'ogicamente realizables}

Sea \(\mathscr H_{\rm HMT}\) la suma numerable de los espacios polacos de
historias HMT tipadas, cada uno provisto de su topolog\'ia de cilindros. Una
\emph{realizaci\'on geneal\'ogicamente admisible} es un subconjunto
\(A\subseteq Q_{\rm HMT}\) para el que existen un lector HMT boreliano
\(F:\mathscr H_{\rm HMT}\to Q_{\rm HMT}\), obtenido por composici\'on de
evaluaciones tipadas de la arquitectura, y una relaci\'on boreliana
\(C\subseteq\mathscr H_{\rm HMT}\times\mathbb N^{\mathbb N}\) tales que
\begin{equation}
 A=\{F(x):\exists z\in\mathbb N^{\mathbb N},\ (x,z)\in C\}.
 \label{eq:hmt-analytic-publication}
\end{equation}
El testigo \(z\) codifica una prolongaci\'on completa, una ruta, una memoria
o un certificado numerable; \(F\) act\'ua sobre una historia ya construida.
Se denota por
\begin{equation}
 \operatorname{Real}_{\rm HMT}(Q_{\rm HMT})
 :=\{A\subseteq Q_{\rm HMT}:A\text{ satisface }
 \eqref{eq:hmt-analytic-publication}\}
 \label{eq:clase-realizaciones-genealogicas}
\end{equation}
la clase de realizaciones de la teor\'ia. Esta definici\'on fija el tipo de
los cuantificadores: un subconjunto es un objeto HMT cuando dispone de una
genealog\'ia expresada por un lector y un testigo de construcci\'on. La
operaci\'on conjunto potencia, considerada sin ese dato, pertenece al reducto
extensional y no a la categor\'ia generadora.

\begin{theorem}[Dicotom\'ia cardinal de las realizaciones HMT]
\label{thm:dicotomia-cardinal-realizaciones-hmt}
Toda realizaci\'on
\(A\in\operatorname{Real}_{\rm HMT}(Q_{\rm HMT})\) satisface
\begin{equation}
 \boxed{|A|\leq\aleph_0\qquad\text{o}\qquad
 |A|=2^{\aleph_0}.}
 \label{eq:continuum-analytic-dichotomy}
\end{equation}
Si
\begin{equation}
 \aleph_{1,\rm HMT}^{\rm gen}
 :=\min\{|A|:A\in\operatorname{Real}_{\rm HMT}(Q_{\rm HMT}),\
 |A|>\aleph_0\},
 \label{eq:hmt-analytic-cardinal-definition}
\end{equation}
entonces
\begin{equation}
 \boxed{\aleph_{1,\rm HMT}^{\rm gen}
 =2^{\aleph_0}=|Q_{\rm HMT}|.}
 \label{eq:hmt-analytic-cardinal}
\end{equation}
\end{theorem}

\begin{proof}
La ecuaci\'on \eqref{eq:hmt-analytic-publication} representa \(A\) como
imagen boreliana de la proyecci\'on de un boreliano de un espacio polaco;
por tanto, \(A\) es anal\'itico. Si \(A\) no es numerable, el teorema del
conjunto perfecto proporciona una inyecci\'on continua
\(2^{\mathbb N}\hookrightarrow A\) \cite{Kechris1995}. La inclusi\'on
\(A\subseteq Q_{\rm HMT}\cong\mathbb R\) proporciona una inyecci\'on en
sentido inverso, y Cantor--Bernstein da
\(|A|=2^{\aleph_0}\). El portador completo pertenece a
\(\operatorname{Real}_{\rm HMT}(Q_{\rm HMT})\): basta tomar la evaluaci\'on
total como lector y la relaci\'on universal como testigo. Por ello, el
m\'inimo de \eqref{eq:hmt-analytic-cardinal-definition} existe y satisface
\eqref{eq:hmt-analytic-cardinal}.
\end{proof}

\begin{corollary}[Inexistencia de un cardinal intermedio admisible]
\label{cor:inexistencia-cardinal-intermedio-hmt}
No existe una realizaci\'on geneal\'ogicamente admisible \(A\) tal que
\begin{equation}
 \aleph_0<|A|<|Q_{\rm HMT}|.
 \label{eq:inexistencia-cardinal-intermedio-hmt}
\end{equation}
Equivalentemente, la sentencia del continuo bien tipada en HMT,
\begin{equation}
 \mathsf{CH}_{\rm HMT}:
 \quad
 \nexists A\in\operatorname{Real}_{\rm HMT}(Q_{\rm HMT})
 \ \bigl(\aleph_0<|A|<|Q_{\rm HMT}|\bigr),
 \label{eq:ch-hmt}
\end{equation}
es un teorema.
\end{corollary}

\begin{remark}[Autonomía de la sentencia cardinal HMT]
\label{rem:autonomia-ch-hmt-autonomo}
La composición APP--TRIT--TPK produce el espacio de historias, su límite
coinductivo y el portador \(Q_{\rm HMT}\). Los lectores tipados definen
después \(\operatorname{Real}_{\rm HMT}(Q_{\rm HMT})\), y el teorema del
conjunto perfecto certifica la dicotomía cardinal de sus imágenes
analíticas. Por tanto, \(\mathsf{CH}_{\rm HMT}\) no presupone
\(\mathsf{CH}_{\rm ext}\), un buen orden de la recta, \(L[h]\), una
aplicación ordinal \(\rho_h\) ni una identificación previa con
\(\aleph_1\). Esas construcciones pertenecen a la comparación extensional
posterior.
\end{remark}

\subsection{Fibras arquimedianas y no fidelidad del funtor de olvido extensional}

El atlas integral \(\mathcal A_{\rm HMT}\) posee la evaluación
\(P_{\rm HMT}\) de \eqref{eq:evaluacion-hmt-cardinal}. Su relación de fibras
contrae exactamente las memorias Hensel brutas que realizan una misma
coordenada. El sistema enriquecido TPK posee, además, un funtor de olvido extensional que
conserva la fase visible y contrae acarreo y memoria. Ambos olvidos se
distinguen por sus dominios.

\begin{proposition}[Obstrucci\'on de factorizaci\'on de los invariantes geneal\'ogicos]
\label{prop:no-fidelidad-evaluacion-arquimediana}
Sea \(I:\mathcal A_{\rm HMT}\to Y\) un invariante geneal\'ogico. Si existen
\(x,y\in\mathcal A_{\rm HMT}\) con
\(P_{\rm HMT}(x)=P_{\rm HMT}(y)\) e \(I(x)\neq I(y)\), entonces no existe
\(F:Q_{\rm HMT}\to Y\) tal que \(I=F\circ P_{\rm HMT}\).
\end{proposition}

\begin{proof}
La igualdad de las evaluaciones implicar\'ia
\(I(x)=F(P_{\rm HMT}(x))=F(P_{\rm HMT}(y))=I(y)\), contradicci\'on.
\end{proof}

Sea \(\rho\) una historia TPK completa y sea
\(\mathcal H_\rho\) la categoría de un objeto con
\begin{equation}
 \operatorname{End}_{\mathcal H_\rho}(\rho)
 =\{\Gamma_9^a:a\in\mathbb N\},
 \qquad
 \Gamma_9^a\circ\Gamma_9^b=\Gamma_9^{a+b}.
 \label{eq:monoide-holonomia-cardinal}
\end{equation}
La graduaci\'on conjunta de acarreo y memoria satisface
\begin{equation}
 \operatorname{gr}_{\rm cm}(\operatorname{id}_\rho)=(0,0),
 \qquad
 \operatorname{gr}_{\rm cm}(\Gamma_9)=(1,1),
 \label{eq:grados-identidad-holonomia}
\end{equation}
y el funtor de olvido extensional
\begin{equation}
 \mathsf V_{\rm ext}:\mathcal H_\rho\longrightarrow\mathbf{Set},
 \qquad
 \mathsf V_{\rm ext}(\rho)=\{\bar\rho\},
 \qquad
 \mathsf V_{\rm ext}(\Gamma_9^a)
 =\operatorname{id}_{\{\bar\rho\}},
 \label{eq:funtor-olvido-arquimediano}
\end{equation}
cumple
\begin{equation}
 \mathsf V_{\rm ext}(\operatorname{id}_\rho)
 =\mathsf V_{\rm ext}(\Gamma_9).
 \label{eq:no-fidelidad-morfismos}
\end{equation}
Por tanto, \(\mathsf V_{\rm ext}\) presenta un defecto de fidelidad: identifica
morfismos con grados distintos de acarreo y memoria. La clase geneal\'ogica
del endomorfismo requiere, además de la fase extensional, su grado de acarreo
y su memoria.

\begin{corollary}[Insuficiencia certificadora del reducto extensional]
\label{cor:insuficiencia-certificadora-extensional-autonomo}
Ningún certificado que distinga la identidad genealógica, el retorno de
fase, el avance de memoria y la holonomía \(\Gamma_9\) puede factorizar
exclusivamente a través de \(\mathsf V_{\rm ext}\). En consecuencia, la
estructura desnuda de pertenencia de ZF o ZFC no constituye por sí sola un
certificado completo del continuo HMT.
\end{corollary}

\begin{proof}
Si \(C=\widetilde C\circ\mathsf V_{\rm ext}\), entonces
\eqref{eq:no-fidelidad-morfismos} fuerza
\[
 C(\operatorname{id}_\rho)
 =\widetilde C\bigl(\mathsf V_{\rm ext}(\operatorname{id}_\rho)\bigr)
 =\widetilde C\bigl(\mathsf V_{\rm ext}(\Gamma_9)\bigr)
 =C(\Gamma_9),
\]
mientras \eqref{eq:grados-identidad-holonomia} distingue ambos
endomorfismos. ZF/ZFC puede codificar la categoría enriquecida una vez
suministrados sus objetos, flechas y grados; el reducto extensional solo no
los recupera.
\end{proof}

Este hecho precisa el defecto fundacional de la reducci\'on extensional. ZFC
puede codificar \(\mathscr G_{\rm HMT}\), sus fibras y sus invariantes una
vez que se le suministran. Sus primitivas de pertenencia, extensionalidad y
conjunto potencia no recuperan, a partir de una coordenada real, la
transici\'on que la produjo, la orientaci\'on conservada, la memoria
transportada ni el estrato dimensional de procedencia. La representabilidad
de una genealog\'ia como conjunto y su determinaci\'on por una ley generadora
son propiedades matem\'aticas distintas.

\subsection{Comparaci\'on exacta con la sentencia extensional de ZFC}

La sentencia \(\mathsf{CH}_{\rm ext}\) de
\eqref{eq:ch-exacta-sobre-portador-hmt} posee la siguiente caracterizaci\'on
ordinal, que se registra exclusivamente para identificar el reducto
extensional.

\begin{theorem}[Criterio ordinal de la sentencia extensional]
\label{thm:criterio-ordinal-ch}
En ZFC son equivalentes:
\begin{equation}
\begin{aligned}
2^{\aleph_0}=\aleph_1
&\Longleftrightarrow Q_{\rm HMT}\hookrightarrow\omega_1\\
&\Longleftrightarrow
\exists\rho:\mathcal A_{\rm HMT}\to\omega_1,\quad
\rho(x)=\rho(y)\Longleftrightarrow
P_{\rm HMT}(x)=P_{\rm HMT}(y)\\
&\Longleftrightarrow
Q_{\rm HMT}\text{ admite un buen orden de tipo }\omega_1.
\end{aligned}
\label{eq:classical-ch-equivalences}
\end{equation}
\end{theorem}

\begin{proof}
Existe siempre una inyecci\'on
\(\omega_1\hookrightarrow Q_{\rm HMT}\cong\mathbb R\). Por
Cantor--Bernstein, la igualdad \(|Q_{\rm HMT}|=\aleph_1\) equivale a la
inyecci\'on inversa. Su composici\'on con \(P_{\rm HMT}\) produce \(\rho\);
la condici\'on sobre las fibras hace descender \(\rho\) a una inyecci\'on del
cociente. Una biyecci\'on con \(\omega_1\) transporta su buen orden y,
rec\'iprocamente, un buen orden de ese tipo produce la biyecci\'on.
\end{proof}

La caracterización ordinal precedente no posee una realización continua o
analítica. Si \(\mathcal A_{\rm HMT}\) es unión numerable de cartas compactas,
no existe una aplicación continua
\(\mathcal A_{\rm HMT}\to\omega_1\), con la topología de orden, que separe
todas las fibras de \(P_{\rm HMT}\): la imagen de cada carta compacta es
acotada y numerable, y una unión numerable de tales imágenes continúa siendo
numerable. Tampoco existe un buen orden analítico de \(\mathbb R\), pues una
relación analítica bien fundada sobre un espacio polaco tiene rango ordinal
numerable. Por consiguiente, los separadores ordinales que aparecen a
continuación son objetos conjuntistas definibles en sus universos
respectivos, no lectores continuos ni analíticos del atlas.

\subsection{Comparación con teorías extensionales del continuo}

La formulación cardinal de Cantor y el primer problema de Hilbert fijan la
pregunta extensional \(2^{\aleph_0}=\aleph_1\). La teoría de conjuntos de
Zermelo--Fraenkel, con o sin el axioma de elección (ZF/ZFC), formaliza esa
pregunta sobre totalidades constituidas mediante la relación de pertenencia.
La jerarquía de constructibilidad de Gödel realiza \(\mathsf{GCH}\), mientras
el método de Cohen produce extensiones de modelos de ZFC en las que la función
del continuo adopta valores diferentes \cite{Godel1940,Cohen1963}. El teorema
de Easton determina, sobre los cardinales regulares, la libertad de la función
potencia compatible con monotonía y cofinalidad \cite{Easton1970}.

El axioma de Martin (MA) actúa como principio de selección para familias de
órdenes parciales; el régimen usual \(\mathsf{MA}+\neg\mathsf{CH}\) constituye
una comparación extensional posterior \cite{MartinSolovay1970}. El axioma de
forzamiento propio (PFA) y Martin's Maximum (MM) fijan en sus respectivos
marcos \(2^{\aleph_0}=\aleph_2\); los desarrollos de
\(\mathsf{MM}^{++}\) y del axioma \((*)\) fortalecen la teoría de
\(H_{\omega_2}\) \cite{Baumgartner1984,ForemanMagidorShelah1988,AsperoSchindler2021}.
Los grandes cardinales proporcionan fuerza de consistencia y principios de
reflexión; el teorema de Lévy--Solovay conserva cardinales medibles bajo
forzamientos pequeños y mantiene disponibles tanto CH como su negación
\cite{LevySolovay1967}. La absolutidad genérica y la \(\Omega\)-lógica
estudian regímenes de invariancia para teorías del continuo
\cite{Woodin2001}. El programa de modelos internos culmina con la conjetura
Ultimate \(L\), mientras la interpretación multiversista organiza la
variación de CH entre universos \cite{Woodin2010,Hamkins2012}.

Cada uno de estos marcos aporta un axioma de selección, un lector de modelos o
un régimen de estabilización del reducto extensional y actúa sobre la función
\(\kappa\mapsto2^\kappa\) una vez fijado el lenguaje conjuntista. El teorema
HMT ocupa una residencia anterior y tipada: construye el portador mediante
el transductor integral, determina el sistema de historias, conserva sus
fibras bajo el TPK y define la clase de realizaciones por lectores
genealógicos. La igualdad
\(\aleph_{1,\rm HMT}^{\rm gen}=2^{\aleph_0}\) es un invariante de esa
categoría generada. El forzamiento, la constructibilidad, los axiomas de
forzamiento, los grandes cardinales, la \(\Omega\)-lógica, Ultimate \(L\) y
el multiverso constituyen comparaciones posteriores de su reducto extensional.
La ley APP--TRIT--TPK produce previamente \(\mathcal A_{\rm HMT}\),
\(Q_{\rm HMT}\) y \(\operatorname{Real}_{\rm HMT}(Q_{\rm HMT})\), y su
teorema cardinal conserva validez con independencia del programa extensional
empleado para representar posteriormente la salida.

\begin{theorem}[Separaci\'on entre independencia extensional y clausura geneal\'ogica]
\label{thm:alcance-godel-cohen-reducto}
G\"odel y Cohen establecen la independencia de
\(\mathsf{CH}_{\rm ext}\) respecto de ZFC. Separadamente,
\(\mathsf{CH}_{\rm HMT}\) es un teorema descriptivo de la categor\'ia de
realizaciones geneal\'ogicamente admisibles. La no fidelidad de
\(\mathsf V_{\rm ext}\) impide identificar ambas sentencias, porque sus
dominios de cuantificaci\'on y sus morfismos no coinciden. En particular:
\begin{enumerate}[label=\textup{(\roman*)}]
 \item \(\mathsf{CH}_{\rm HMT}\) es consecuencia de la dicotom\'ia cardinal
 de las realizaciones geneal\'ogicamente admisibles;
 \item \(\mathsf{CH}_{\rm ext}\) cuantifica sobre
 \(\mathcal P(Q_{\rm HMT})\), incluso sobre subconjuntos sin lector ni
 testigo generador;
 \item la no fidelidad de \(\mathsf V_{\rm ext}\) impide identificar ambos
 dominios de cuantificaci\'on mediante el mero olvido de estructura;
 \item ning\'un resultado de independencia para
 \(\mathsf{CH}_{\rm ext}\) contradice
 \eqref{eq:inexistencia-cardinal-intermedio-hmt}.
\end{enumerate}
La no fidelidad del reducto extensional no se identifica con una causa de
la independencia l\'ogica: una afirmaci\'on de esa naturaleza exigir\'ia una
interpretaci\'on entre modelos compatible con la relaci\'on de forzamiento.
\end{theorem}

\begin{proof}
El primer apartado es el
\cref{cor:inexistencia-cardinal-intermedio-hmt}. El segundo sigue de
\eqref{eq:ch-exacta-sobre-portador-hmt}. El tercero es consecuencia de
\eqref{eq:no-fidelidad-morfismos}: el reducto extensional identifica
morfismos e invariantes geneal\'ogicamente distintos y, por tanto, no puede
recuperar el tipo de realizaci\'on a partir de la sola pertenencia. G\"odel y
Cohen prueban, mediante constructibilidad y forzamiento, la independencia de
la sentencia extensional relativa a ZFC \cite{Godel1940,Cohen1963}; esa
sentencia posee un dominio de cuantificaci\'on distinto del fijado por
\eqref{eq:clase-realizaciones-genealogicas}. La diferencia de dominio y la
no fidelidad impiden transferir por simple olvido una conclusi\'on entre las
dos sentencias. De esta separaci\'on se sigue el cuarto apartado, sin atribuir
al funtor de olvido el origen metamatem\'atico de la independencia.
\end{proof}

\begin{proposition}[Cambio de base por forzamiento sobre objetos heredados]
\label{prop:cuadrado-forcing-heredado-hmt-autonomo}
Sea \(M\) un modelo transitivo de ZFC que contiene el código del generador
APP--TRIT--TPK, sea \(\mathbb P\in M\), sea \(G\subseteq\mathbb P\)
\(M\)-genérico y sea \(N=M[G]\). Denotemos por
\(\mathcal H_{\rm HMT}^{M}\) la categoría de historias y morfismos HMT
pertenecientes a \(M\), y por
\(\left.\mathcal H_{\rm HMT}^{N}\right|_{M}\) la subcategoría de \(N\)
formada por esos mismos objetos y flechas heredados. Escribimos
\(\mathsf V_{\rm ext}^{M}\) y \(\mathsf V_{\rm ext}^{N}\) para los
funtores de olvido de \eqref{eq:funtor-olvido-arquimediano} interpretados
en cada modelo. Las inclusiones
canónicas \(\iota_G^{\mathcal H}\) y \(\iota_G^{\rm set}\) forman el
cuadrado
\begin{equation}
\begin{array}{ccc}
 \mathcal H_{\rm HMT}^{M}
 &\xrightarrow{\ \iota_G^{\mathcal H}\ }&
 \left.\mathcal H_{\rm HMT}^{N}\right|_{M}\\[3pt]
 {\scriptstyle\mathsf V_{\rm ext}^{M}}\!\downarrow
 &&\downarrow\!{\scriptstyle\mathsf V_{\rm ext}^{N}}\\[3pt]
 \mathbf{Set}^{M}
 &\xrightarrow{\ \iota_G^{\rm set}\ }&
 \left.\mathbf{Set}^{N}\right|_{M},
\end{array}
\qquad
 \mathsf V_{\rm ext}^{N}\circ\iota_G^{\mathcal H}
 =\iota_G^{\rm set}\circ\mathsf V_{\rm ext}^{M}.
\label{eq:cuadrado-forcing-heredado-hmt-autonomo}
\end{equation}
La compatibilidad se limita a objetos heredados. No afirma que toda
historia, rama, código o subconjunto nuevo de \(N\) proceda de \(M\), ni
establece una equivalencia o una naturalidad fuerte entre las categorías
completas de ambos modelos.
\end{proposition}

\begin{proof}
La inclusión genérica conserva literalmente cada historia y cada morfismo
codificado en \(M\). Aplicar el olvido extensional antes o después de la
inclusión produce el mismo conjunto subyacente y la misma aplicación
subyacente en \(N\), lo que demuestra
\eqref{eq:cuadrado-forcing-heredado-hmt-autonomo}. El argumento no alcanza
los objetos añadidos por el forzamiento y no implica esencial
sobreyectividad.
\end{proof}

La \cref{prop:cuadrado-forcing-heredado-hmt-autonomo} establece la
compatibilidad mínima del cambio de base sobre objetos heredados. Para cada
\(W\in\{M,N\}\), las realizaciones HMT construidas en
\(W\) satisfacen \(W\models\mathsf{CH}_{\rm HMT}\) por la aplicación
interna del \cref{thm:dicotomia-cardinal-realizaciones-hmt}; las ramas nuevas
de \(N\), si las hay, quedan cubiertas por esa demostración en \(N\), no por
una prolongación del cuadrado. El forzamiento puede variar
\(\mathsf{CH}_{\rm ext}\), porque ésta cuantifica sobre todo
\(\mathcal P(Q_{\rm HMT})\), sin alterar el teorema cardinal de la categoría
genealógicamente admisible.

\subsection{Clausura ordinal en el envolvente constructible del generador}

Sea \(h\in2^{\mathbb N}\) el código de una presentación numerable del
lenguaje, las tablas finitas, las transiciones y los lectores
APP--TRIT--TPK. El código no contiene un buen orden previo de la recta, la
relación de verdad de un universo, la sentencia \(\mathsf{CH}\) ni valores
objetivo. Sea
\begin{equation}
 M_h:=L[h].
 \label{eq:envolvente-constructible-generador}
\end{equation}
El parámetro \(h\) conserva el generador como dato interno del modelo, y el
buen orden canónico de la jerarquía constructible relativa proporciona una
coordenada ordinal definible.

\begin{theorem}[Clausura cardinal en \(L\lbrack h\rbrack\)]
\label{thm:ch-envolvente-constructible-generador}
El universo \(M_h=L[h]\) satisface
\begin{equation}
 M_h\models2^{\aleph_0}=\aleph_1.
 \label{eq:ch-envolvente-constructible-generador}
\end{equation}
Existe en \(M_h\) una biyección definible
\begin{equation}
 b_h:\omega_1^{M_h}\xrightarrow{\;\sim\;}\mathbb R^{M_h}.
 \label{eq:biyeccion-envolvente-constructible}
\end{equation}
Si
\begin{equation}
 P_h:=\left.P_{\rm HMT}\right|_{\mathcal A_{\rm HMT}\cap M_h}:
 \mathcal A_{\rm HMT}\cap M_h\longrightarrow\mathbb R^{M_h},
 \label{eq:evaluacion-envolvente-constructible}
\end{equation}
entonces
\begin{equation}
 \rho_h:=b_h^{-1}\circ P_h:
 \mathcal A_{\rm HMT}\cap M_h\longrightarrow\omega_1^{M_h}
 \label{eq:separador-envolvente-constructible}
\end{equation}
satisface
\begin{equation}
 \rho_h(x)=\rho_h(y)
 \quad\Longleftrightarrow\quad
 P_h(x)=P_h(y).
 \label{eq:fibras-envolvente-constructible}
\end{equation}
\end{theorem}

\begin{proof}
Escribamos \(M=L[h]\). Para cada \(x\in\mathbb R^M\), un casco de Skolem
numerable de algún nivel \(L_\eta[h]\) que contenga \(h,x,\omega\) posee un
colapso transitivo que fija \(h\) y \(x\). Por condensación relativizada, el
colapso es \(L_{\bar\theta}[h]\) para algún
\(\bar\theta<\omega_1^M\). Por consiguiente, todo real de \(M\) aparece
antes de \(\omega_1^M\), y cada nivel de primera aparición es numerable en
\(M\) \cite{Jech2003}. Se sigue que
\[
 |\mathbb R^{L[h]}|\leq\aleph_1^{L[h]}.
\]
La recta del modelo es no numerable, y la minimalidad de
\(\aleph_1^{L[h]}\) da la desigualdad inversa. El buen orden canónico de
\(L[h]\) selecciona una biyección definible \(b_h\). La evaluación \(P_h\)
es definible a partir de \(h\) y toma valores en \(\mathbb R^{M_h}\); la
biyectividad de \(b_h\) demuestra
\eqref{eq:fibras-envolvente-constructible}.
\end{proof}

El separador \(\rho_h\) clasifica todas las fibras de evaluación cuyos
datos pertenecen a \(M_h\). Para un universo ambiente
\(V\supseteq M_h\), la existencia de algún separador
\begin{equation}
 \rho_V:\mathcal A_{\rm HMT}^V\longrightarrow\omega_1^V,
 \qquad
 \rho_V(x)=\rho_V(y)\Longleftrightarrow P_V(x)=P_V(y),
 \label{eq:separador-ordinal-ambiente}
\end{equation}
equivale a una inyección
\(Q_{\rm HMT}^V\hookrightarrow\omega_1^V\) y, por
\cref{thm:criterio-ordinal-ch}, a
\(V\models2^{\aleph_0}=\aleph_1\). No se afirma que \(\rho_V\) prolongue
puntualmente la función \(\rho_h\): la comparación de ambos separadores
puede exigir una reindexación de \(\omega_1^{M_h}\) en \(\omega_1^V\).
Quedan así separados tres resultados:
la dicotomía cardinal de las realizaciones genealógicas, la equivalencia
ordinal exacta de la sentencia extensional y la clausura efectiva de todas
las fibras internas de \(L[h]\).

El estatuto l\'ogico del resultado es
\(\mathsf{CH}_{\rm HMT}\), mientras la igualdad ambiental
\(V\models2^{\aleph_0}=\aleph_1\) constituye la sentencia distinta
\(\mathsf{CH}_{\rm ext}\). La teor\'ia HMT construye el continuo, determina la clase de sus
realizaciones matem\'aticamente admisibles y demuestra que en esa clase no
existe cardinal intermedio. La independencia de la sentencia extensional
relativa a ZFC y la clausura cardinal de la categor\'ia geneal\'ogica son
resultados de dominios distintos; la no fidelidad del olvido impide
transferirlos por identificaci\'on de sus cuantificadores.

\subsection{Clausura cardinal en la categoría genealógica}

\begin{theorem}[Cierre cardinal del continuo generado]
\label{thm:resolucion-genealogica-continuo}
La arquitectura APP--TRIT--TPK determina, a partir de sus datos finitos y de
sus leyes tipadas de prolongaci\'on:
\begin{enumerate}[label=\textup{(\roman*)}]
 \item la estructura discreta conjunta
 \(\mathfrak C_{\rm disc}\) como l\'imite inverso de historias enriquecidas;
 \item el portador lineal
 \(Q_{\rm HMT}\cong\mathbb R_{\rm HMT}\cong\mathbb R\) y, para todo
 \(d\geq1\), los portadores dimensionales
 \(Q_{\rm HMT}^{\,d}\cong\mathbb R^d\);
 \item las fibras de evaluaci\'on que conservan hoja, orientaci\'on, fase,
 acarreo, frontera, ruta y memoria aunque la coordenada real las contraiga;
 \item la clase \(\operatorname{Real}_{\rm HMT}(Q_{\rm HMT})\) de
 realizaciones geneal\'ogicamente admisibles;
 \item la igualdad
 \(\aleph_{1,\rm HMT}^{\rm gen}=2^{\aleph_0}\) y, por consiguiente, la
 inexistencia de un cardinal intermedio admisible;
 \item para el código completo \(h\) del generador, la clausura
 constructible relativa \(M_h=L[h]\), la igualdad
 \(M_h\models2^{\aleph_0}=\aleph_1\) y el separador definible
 \(\rho_h=b_h^{-1}\circ P_h\) de todas las fibras internas;
 \item la no fidelidad del reducto extensional y la separaci\'on formal entre
 la independencia de G\"odel--Cohen para \(\mathsf{CH}_{\rm ext}\) y la
 clausura cardinal demostrada para \(\mathsf{CH}_{\rm HMT}\).
\end{enumerate}
Por tanto, el problema del continuo queda resuelto en la teor\'ia HMT: el
continuo lineal y dimensional se construye, su portador se realiza como
cociente y su espectro cardinal geneal\'ogicamente admisible contiene
exactamente el estrato numerable y el estrato continuo, sin cardinal
intermedio.
\end{theorem}

\begin{proof}
Los apartados \textup{(i)}--\textup{(iii)} se siguen de la construcci\'on
proyectiva de la \cref{sec:construccion-proyectiva-recta}, del
\cref{thm:cociente-arquimediano-atlas-nuclear} y de su producto dimensional.
El apartado \textup{(iv)} es la
\cref{eq:clase-realizaciones-genealogicas}; el apartado \textup{(v)} es el
\cref{thm:dicotomia-cardinal-realizaciones-hmt}; el apartado \textup{(vi)}
es el \cref{thm:ch-envolvente-constructible-generador}; y el apartado
\textup{(vii)} resulta de la
\cref{prop:no-fidelidad-evaluacion-arquimediana,thm:alcance-godel-cohen-reducto}.
\end{proof}
